\clearpage
\begin{landscape}
\thispagestyle{plain}

\begin{figure}[p]
\centering
\vspace*{-0.4cm}
{\Large\bfseries\color{MinovaDark} UML-Komponentendiagramm\par}
\vspace{0.2cm}
{\small\color{MinovaDark}
Komponenten, Abhängigkeiten und Benutzer der digitalen Lieferavisierung\par}
\vspace{0.45cm}

\resizebox{0.98\linewidth}{!}{%
\begin{tikzpicture}[
  font=\small,
  internal/.style={
    draw=MinovaDark,
    line width=1.1pt,
    minimum width=4.85cm,
    minimum height=1.55cm,
    align=center,
    fill=SoftBlue,
    inner sep=6pt
  },
  external/.style={
    draw=MinovaDark,
    line width=1.1pt,
    minimum width=4.25cm,
    minimum height=1.5cm,
    align=center,
    fill=SoftGray,
    inner sep=6pt
  },
  note/.style={
    draw=MinovaDark,
    line width=0.8pt,
    minimum width=4.35cm,
    minimum height=1.3cm,
    align=center,
    fill=SoftGray,
    inner sep=6pt
  },
  dependency/.style={-{Latex[open,length=2.5mm,width=1.8mm]},line width=0.9pt,dashed,draw=MinovaDark},
  association/.style={line width=0.9pt,draw=MinovaDark},
  noteLink/.style={line width=0.8pt,dashed,draw=MinovaDark},
  edgeLabel/.style={font=\small,align=center,fill=white,inner sep=2.5pt,text=MinovaDark},
  actorLabel/.style={font=\bfseries,align=center,text=MinovaDark}
]
  % Externe Systeme und Akteure
  \node[external] (dispo) at (3.15,0.7)
    {\textbf{DISPO}\\Führendes Planungssystem};

  % UML-Notizen: fachliche Hinweise, keine Komponenten
  \node[note,minimum width=4.8cm] (integrationNote) at (3.15,3.05)
    {DISPO-Integration ist noch\\fachlich abzustimmen};
  \node[note] (deliveryNote) at (21.45,0.7)
    {Benachrichtigung über\\E-Mail, SMS oder WhatsApp};

  % Fachliche Komponenten innerhalb der Systemgrenze
  \node[internal] (service) at (12.3,0.7)
    {\textbf{Avisierungsservice}\\Versand und Rückmeldungen};
  \node[internal] (dashboard) at (8.65,-3.85)
    {\textbf{Avisierungsdashboard}\\Interne Disponentenansicht};
  \node[internal] (page) at (15.95,-3.85)
    {\textbf{Bestätigungsseite}\\Öffentliche Kundenansicht};

  % UML-Komponentensymbole
  \foreach \componentNode in {dispo,service,dashboard,page}{
    \draw[line width=0.65pt,draw=MinovaDark,fill=white]
      ([xshift=-0.65cm,yshift=-0.22cm]\componentNode.north east)
      rectangle ++(0.43cm,-0.31cm);
    \draw[line width=0.65pt,draw=MinovaDark,fill=white]
      ([xshift=-0.76cm,yshift=-0.27cm]\componentNode.north east)
      rectangle ++(0.18cm,-0.08cm);
    \draw[line width=0.65pt,draw=MinovaDark,fill=white]
      ([xshift=-0.76cm,yshift=-0.40cm]\componentNode.north east)
      rectangle ++(0.18cm,-0.08cm);
  }

  % Gefaltete Ecken der UML-Notizen
  \foreach \noteNode in {integrationNote,deliveryNote}{
    \fill[white]
      ([xshift=-0.36cm]\noteNode.north east) --
      (\noteNode.north east) --
      ([yshift=-0.36cm]\noteNode.north east) -- cycle;
    \draw[line width=0.8pt,draw=MinovaDark]
      ([xshift=-0.36cm]\noteNode.north east) --
      ([xshift=-0.36cm,yshift=-0.36cm]\noteNode.north east) --
      ([yshift=-0.36cm]\noteNode.north east);
  }

  % Akteur: Disponent
  \coordinate (dispatcherHead) at (3.15,-3.35);
  \draw[line width=0.9pt,draw=MinovaDark] (dispatcherHead) circle (0.23cm);
  \draw[line width=0.9pt,draw=MinovaDark]
    (3.15,-3.58) -- (3.15,-4.25)
    (2.70,-3.82) -- (3.60,-3.82)
    (3.15,-4.25) -- (2.78,-4.76)
    (3.15,-4.25) -- (3.52,-4.76);
  \node[actorLabel,anchor=north] (dispatcher) at (3.15,-4.87) {Disponent};

  % Akteur: Kunde
  \coordinate (customerHead) at (21.45,-3.35);
  \draw[line width=0.9pt,draw=MinovaDark] (customerHead) circle (0.23cm);
  \draw[line width=0.9pt,draw=MinovaDark]
    (21.45,-3.58) -- (21.45,-4.25)
    (21.00,-3.82) -- (21.90,-3.82)
    (21.45,-4.25) -- (21.08,-4.76)
    (21.45,-4.25) -- (21.82,-4.76);
  \node[actorLabel,anchor=north] (customer) at (21.45,-4.87) {Kunde};

  % Das Gesamtsystem wird als zusammengesetzte UML-Komponente dargestellt
  \begin{scope}[on background layer]
    \node[
      draw=MinovaDark,
      line width=1.2pt,
      fill=SoftBlue!22,
      fit=(service)(dashboard)(page),
      inner xsep=0.72cm,
      inner ysep=1.1cm,
      label={[font=\bfseries\large,text=MinovaDark,fill=white,inner sep=3pt]above:
        System zur digitalen Lieferavisierung}
    ] (systemBoundary) {};
  \end{scope}

  % Komponentensymbol des Gesamtsystems
  \draw[line width=0.7pt,draw=MinovaDark,fill=white]
    ([xshift=-0.72cm,yshift=-0.22cm]systemBoundary.north east)
    rectangle ++(0.48cm,-0.34cm);
  \draw[line width=0.7pt,draw=MinovaDark,fill=white]
    ([xshift=-0.84cm,yshift=-0.28cm]systemBoundary.north east)
    rectangle ++(0.19cm,-0.08cm);
  \draw[line width=0.7pt,draw=MinovaDark,fill=white]
    ([xshift=-0.84cm,yshift=-0.42cm]systemBoundary.north east)
    rectangle ++(0.19cm,-0.08cm);

  % UML-Abhängigkeiten zwischen Komponenten
  \draw[dependency]
    ([yshift=0.25cm]dispo.east) --
    node[edgeLabel,above] {Avisierung starten\\Liefer- und Tourdaten}
    ([yshift=0.25cm]service.west);
  \draw[dependency]
    ([yshift=-0.28cm]service.west) --
    node[edgeLabel,below] {Ergebnis-Callback}
    ([yshift=-0.28cm]dispo.east);
  \draw[dependency,rounded corners=4pt]
    (page.north) -- (15.95,-1.62) --
    node[edgeLabel,above] {Lieferdaten und Kundenantwort}
    (14.55,-1.62) -- (service.south east);
  \draw[dependency,rounded corners=4pt]
    (dashboard.north) -- (8.65,-1.62) --
    node[edgeLabel,above] {Fälle, Kommentare\\und Aktionen}
    (10.05,-1.62) -- (service.south west);

  % UML-Assoziationen der Benutzer mit ihren Oberflächen
  \draw[association] (3.15,-3.05) --
    node[edgeLabel,right] {Planung}
    (dispo.south);
  \draw[association] (3.60,-4.05) --
    node[edgeLabel,above] {Fälle bearbeiten}
    (dashboard.west);
  \draw[association] (21.0,-4.05) --
    node[edgeLabel,above] {Bestätigen, ablehnen\\oder kommentieren}
    (page.east);

  % Verbindungen von UML-Notizen zu den beschriebenen Elementen
  \draw[noteLink] (integrationNote.south) -- (dispo.north);
  \draw[noteLink] (service.east) -- (deliveryNote.west);
\end{tikzpicture}%
}

\caption{UML-Komponentendiagramm der digitalen Lieferavisierung}
\label{fig:uml-component-overview}
\end{figure}

\end{landscape}
\clearpage
